% Punkte und Strecken im Koordinatensystem · Dreiecke nachweisen · Prüfungs-Fokus Abitur GK – bank/punkte-und-strecken-im-koordinatensystem/e3.jsonl (zusammenbau v0.6, Prüfungs-Fokus)
\einheitenkopf[e3]{Einheit 3 · Dreiecke nachweisen}
\zweigzeile{Abi GK ×7}
\setcounter{aufgabe}{0}

% Anlauf (ohne Prüfkennung)
\begin{aufgabe}{Dreiecke nachweisen – Anlauf}
\begin{teile}
\teil Parallelogramm $ABCD$: Welcher Vektor ist gleich $\overrightarrow{AB}$? \\ \kreuz{$\overrightarrow{CD}$} \\ \kreuz{$\overrightarrow{DC}$} \\ \kreuz{$\overrightarrow{BC}$}
\end{teile}
\end{aufgabe}

\begin{aufgabe}{Dreiecke nachweisen – Prüfungsaufgaben}
\begin{teile}
\teil Zeige: Das Dreieck $A(2 | 0 | 1)$, $B(6 | 2 | 1)$, $C(4 | 1 | 5)$ ist gleichschenklig. Nenne Basis und Schenkel. \hfill (A23)
\teil Zeige: Das Dreieck $A(0 | 2 | 3)$, $B(4 | 2 | 1)$, $C(2 | 5 | 2)$ ist gleichschenklig. Nenne Basis und Schenkel. \hfill (A23)
\teil Zeige: Das Dreieck $P(1 | 4 | 0)$, $Q(5 | 2 | 0)$, $R(5 | 7 | 2)$ ist gleichschenklig. Nenne Basis und Schenkel. \hfill (A23)
\teil Zeige: Das Dreieck $A(2 | 1 | 4)$, $B(4 | 5 | 2)$, $C(5 | 3 | 5)$ ist gleichschenklig. Nenne Basis und Schenkel. \hfill (A23)
\teil Zeige: Das Dreieck $A(0 | 3 | 1)$, $B(4 | 1 | 5)$, $C(4 | 6 | 3)$ ist gleichschenklig. Nenne Basis und Schenkel. \hfill (A23)
\teil Prüfe, ob das Dreieck $A(4 | 2 | 2)$, $B(2 | 4 | 2)$, $C(2 | 2 | 4)$ gleichseitig ist. \hfill (A23)
\teil Prüfe, ob das Dreieck $A(5 | 2 | 2)$, $B(2 | 5 | 2)$, $C(2 | 2 | 6)$ gleichseitig ist. \hfill (A20)
\teil Prüfe, ob das Dreieck $A(7 | 3 | 1)$, $B(3 | 7 | 1)$, $C(3 | 3 | 5)$ gleichseitig ist. \hfill (A23)
\teil $A(1 | 2 | 0)$, $B(5 | 6 | 0)$ und $C(3 | 4 | p)$ mit $p > 0$. Zeige: Das Dreieck $ABC$ ist für jedes $p$ gleichschenklig mit der Basis $AB$. \hfill (A22)
\teil $A(3 | 0 | 3)$, $B(3 | 6 | 1)$ und $C(p | 3 | 2)$. Zeige: Das Dreieck $ABC$ ist für jedes $p \ne 3$ gleichschenklig mit der Basis $AB$. \hfill (A22)
\teil $A(-1 | 2 | 2)$, $B(3 | 2 | -2)$ und $C(1 | p | 0)$. Zeige: Das Dreieck $ABC$ ist für jedes $p$ gleichschenklig. Welche Seite ist die Basis? \hfill (A22)
\teil Ein Quader hat die Ecken $A(9 | 0 | 0)$, $B(9 | 9 | 0)$ und $C(0 | 9 | 0)$ in der $x_1x_2$-Ebene. Begründe: Das Dreieck $ABC$ ist rechtwinklig und gleichschenklig. Flächeninhalt? \hfill (A18)
\end{teile}
\end{aufgabe}

\begin{aufgabe}{Dreiecke nachweisen – Prüfungsaufgaben – weiter}
\begin{teile}
\teil Grundfläche $A(6 | 1 | 0)$, $B(1 | 4 | 0)$, $C(1 | -2 | 0)$; Deckkante $E(1 | 3 | 2)$, $F(1 | -1 | 2)$. Zeige: $ABC$ ist gleichschenklig. Begründe: $EF$ ist parallel zur Grundfläche. \hfill (A23)
\teil Grundfläche $A(6 | 4 | 3)$, $B(2 | 1 | 3)$, $C(2 | 7 | 3)$; Deckkante $G(1 | 3 | 5)$, $H(1 | 5 | 5)$. Zeige: $ABC$ ist gleichschenklig. Begründe: $GH$ ist parallel zur Grundfläche. \hfill (A23)
\teil Grundfläche $A(5 | 6 | 0)$, $B(3 | 1 | 0)$, $C(7 | 1 | 0)$; Deckkante $P(4 | 2 | 3)$, $Q(6 | 2 | 3)$. Zeige: $ABC$ ist gleichschenklig. Begründe: $PQ$ ist parallel zur Grundfläche. \hfill (A23)
\teil $D(8 | 0 | 2)$, $E(0 | 6 | 2)$, $F(0 | 0 | 2)$ und $M(4 | 3 | 2)$. Begründe: $D$, $E$ und $F$ liegen auf einem Kreis mit dem Mittelpunkt $M$. \hfill (A24)
\teil Begründe: $D(2 | 2 | 5)$, $E(6 | 6 | 1)$ und $F(6 | 2 | 5)$ liegen auf einem Kreis, dessen Mittelpunkt die Mitte von $DE$ ist. \hfill (A24)
\teil Begründe: $P(-1 | 3 | 0)$, $Q(5 | 3 | 0)$ und $R(2 | 3 | 3)$ liegen auf einem Kreis, dessen Mittelpunkt die Mitte von $PQ$ ist. \hfill (A24)
\teil In der $x_1x_2$-Ebene liegen $A(0 | 0)$, $B(6 | 0)$ und $C(3 | 9)$ auf einem Kreis. Koordinaten seines Mittelpunkts $M$? \hfill (A19) \\ M( \leerfeld | \leerfeld )
\teil In der $x_1x_2$-Ebene liegen $A(-2 | 0)$, $B(6 | 0)$ und $C(2 | 8)$ auf einem Kreis. Koordinaten seines Mittelpunkts $M$? \hfill (A19) \\ M( \leerfeld | \leerfeld )
\teil In der $x_1x_2$-Ebene liegen $A(1 | 1)$, $B(7 | 1)$ und $C(4 | 10)$ auf einem Kreis. Koordinaten seines Mittelpunkts $M$? \hfill (A19) \\ M( \leerfeld | \leerfeld )
\end{teile}
\end{aufgabe}

\begin{aufgabe}{Dreiecke nachweisen – Prüfungsaufgaben – weiter}
\begin{teile}
\teil Dreieck $A(2 | 1 | 2)$, $B(4 | 5 | 4)$, $C(5 | 3 | 1)$ mit $|CA| = |CB|$. Gesucht ist ein Punkt $D \ne C$ mit $|DA| = |DB|$ und $|DA| = |CA|$. Koordinaten von $D$? \hfill (A23) \\ D( \leerfeld | \leerfeld | \leerfeld )
\teil Dreieck $A(0 | 2 | 3)$, $B(4 | 2 | -1)$, $C(4 | 3 | 3)$ mit $|CA| = |CB|$. Gesucht ist ein Punkt $D \ne C$ mit $|DA| = |DB|$ und $|DA| = |CA|$. Koordinaten von $D$? \hfill (A23) \\ D( \leerfeld | \leerfeld | \leerfeld )
\teil $A(3 | 1 | 2)$ und $M(1 | 3 | 1)$. Gib zwei Punkte $B$ und $C$ an, sodass $A$, $B$ und $C$ denselben Abstand zu $M$ haben und ein gleichschenkliges Dreieck bilden. \hfill (A24)
\teil $A(2 | 0 | 5)$ und $M(4 | 2 | 4)$. Gib zwei Punkte $B$ und $C$ an, sodass $A$, $B$ und $C$ denselben Abstand zu $M$ haben und ein gleichschenkliges Dreieck bilden. \hfill (A24)
\end{teile}
\end{aufgabe}

\medskip
%% Merkkasten Einheit 3: 6 Zeilen, Vorgabe höchstens fünf (3.1) – ungekürzt
\uebersichtskasten{Seiten als Beträge: die drei Seitenlängen entscheiden – gleichschenklig heißt zwei gleich lange Schenkel (die Basis ist die dritte Seite), gleichseitig alle drei gleich; erst benennen, welche Seiten Schenkel sein sollen, dann rechnen. \\ A(1 | 1 | 0), B(5 | 1 | 0), C(3 | 4 | 2): |CA| = |CB| = √(4 + 9 + 4) = √17 → gleichschenklig mit Basis AB. \\ Rechter Winkel: am behaupteten Eckpunkt das Skalarprodukt der beiden anliegenden Seitenvektoren prüfen – null heißt rechtwinklig; bei achsenparallelen Kanten genügt der Blick auf die Koordinaten. \\ Satz des Thales: hat das Dreieck bei C den rechten Winkel, liegt C auf dem Kreis um den Mittelpunkt von AB durch A und B – drei Punkte liegen auf einem Kreis um M, wenn M Mittelpunkt einer Verbindungsstrecke ist und der dritte Punkt dort den rechten Winkel bildet. \\ Umkreismittelpunkt: gleich weit von allen Ecken – in der Ebene eine Koordinate aus der Symmetrie oder Mittelsenkrechten, die andere aus der Gleichheit zweier Abstände. \\ Mit Parameter: die unbekannte Koordinate ansetzen, die Bedingung (gleiche Schenkel, vorgegebener Umfang) als Gleichung mit Wurzeln aufschreiben – erst die Wurzel isolieren, dann quadrieren.}
